%قبل از شروع مسئله ترجیحاً چیزی ننویسید. اما اگر خواستید بنویسید، در قالب لیستی (فقرات و شمارش) بنویسید.
%برای اضافه کردن نمره، بعدی از شروع‌مسئله، [x] را اضافه کنید و به جای x  نمره مورد نظر را قرار دهید. برای زیرمسئله‌ها هم به همین صورت است.
%\شروع‌مسئله[10]
\شروع‌مسئله [17]
%\متن‌سیاه{سوال عملی اول:}\\
%\متن‌سیاه{دراپ‌تست}\\
در یک کفاشی، \متن‌لاتین{n} سایز مختلف کفش وجود دارد. مردی به داخل مغازه آمده و از مغازه‌دار می‌خواهد که یک کفش متناسب با اندازه‌ی پای او پیدا کند(می‌دانیم دقیقا یکی از کفش های مغازه برای مشتری مناسب است). مغازه‌دار می‌تواند در 
$O(1)$
سایز دو کفش را مقایسه کند و بفهمد که کدام بزرگتر است. اما برای این که بفهمد پای مشتری بزرگ تر است یا یک کفش یا که این دو برابرند، نیاز دارد که کفش را پای مشتری بکند. مشتری از این کار خوشش نمی‌آید و هر بار که بخواهد کفشی را بپوشد، \متن‌لاتین{k} ثانیه لفتش می‌دهد. مغازه‌دار هم که حوصله ندارد می‌خواهد سریع‌تر این کار را انجام دهد. راهی از اردر
$O(n + k\log(n))$
به مشتری ارائه دهید.
\solution{
پاسخ. برای این کار از جستجوی دودویی استفاده می‌کنیم. برای این کار ابتدا تمام کفش ها را در نظر میگیریم. در 
$O(n)$
می‌توانیم میانه‌ی کفش ها از لحاظ اندازه را بدست آوریم. سپس کفش ها را به دو تکه‌ی بزرگتر از میانه و کوچکتر از میانه در 
$O(n)$
تقسیم می‌کنیم. حال امتحان می‌کنیم تا بفهمیم سایز کفش میانه و پای مشتری چه نسبتی به هم دارند. در صورت اندازه بودن کار تمام است. در غیر این صورت یا کفش مورد نظر در نیمه‌ی بزرگتر کفش هاست یا در نیمه‌ی کوچکتر کفش ها. سپس روی کفش های باقی مانده همین الگوریتم را تکرار می‌کنیم. 

حال این راه را تحلیل می‌کنیم. حداکثر
$\log(n)$
بار یک کفش را با پای مشتری مقایسه کرده ایم. زیرا هر بار که این کار را کنیم و کار تمام نشود، مجموعه کفش هایی که جواب در آنهاست نصف می‌شود و این کار حداکثر 
$\log(n)$
بار اتفاق می‌افتد. پس کل مقایسه‌های کفش‌ها و پای مشتری، از 
$O(k\log(n))$
است. 
حال پیدا کردن میانه‌ها را تحلیل می‌کنیم. فرض کنید پیدا کردن میانه‌ی 
$n$
عدد حداکثر به 
$C*n$
عملیات نیاز دارد که 
$C$
یک ثابت است. حال مجموع هزینه‌های پیدا کردن میانه‌ها را حساب می‌کنیم.
$$Time = C*n+C*\frac{n}{2}+C*\frac{n}{4}+...+C*\frac{n}{2^{\log(n)}} = C*n(1+\frac{1}{2}+...+\frac{1}{2^{\log(n)}}) < C * n * 2 < C'*n$$
$$Time = C'*n$$
پس می‌فهمیم هزینه های یافتن کل میانه‌ها نیز از 
$O(n)$
است. پس اردر کل الگوریتم برابر است با:
$O(n + k\log(n))$.

}
%\متن‌سیاه{ورودی:}\\
%در خط اول ورودی، ابتدا k و سپس m داده می‌شود.\\
%\متن‌سیاه{خروجی:}\\
%در تنها سطر خروجی، تعداد آیفون‌های مورد نیاز می‌آید. اگر در هیچ صورتی چنین کاری ممکن نبود ۱- چاپ می‌شود.\\
%\متن‌سیاه{محدودیت‌ها:}\\
%$1\leq k\leq 10^4$

%$m\leq10^9$
%\\
%\textbf{\متن‌سیاه{ورودی و خروجی نمونه}\\
%% تست اول
%\begin{latin}
%	\begin{tabularx}{\textwidth}{|X | X|}
%		\hline
%		\textbf{stdin} &
%		\textbf{stdout}\\
%		\hline
%		100 100&
%		1\\
%		\hline
%	
%	\end{tabularx}
%\end{latin}
%
%% تست دوم
%
%\begin{latin}
%	\begin{tabularx}{\textwidth}{|X | X|}
%		\hline
%		\textbf{stdin} &
%		\textbf{stdout}\\
%		\hline
%		340 20&
%		3\\
%		\hline
%		
%	\end{tabularx}
%\end{latin}}
% صورت اصلی مسئله را در این‌جا بنویسید.
%اگر مسئله چند قسمت داشت به این صورت عمل کنید:
%\شروع{زیرمسائل}
%\زیرمسئله[0] % مسئلۀ ۱آ
%\زیرمسئله % مسئلۀ ۱ب
%\زیرمسئله % مسئلۀ ۱پ
%\زیرمسئله % مسئلۀ ۱ت
%\پایان{زیرمسائل}





%%%%%%%%%%%%%%%%%%%%%%%%%%%%%%%%%%%%%%%%%%%%%%%%%%%%%
% برخی دستورات مفید برای تمرین‌ها. %
%%%%%%%%%%%%%%%%%%%%%%%%%%%%%%%%%%%%%%%%%%%%%%%%%%%%%

% فرمول‌های نمونه:
%$f(x) = \Theta(x)$
%$T(x, y) \leq \log(x) + 2^y + \binom{2n}{n}$
% {\tt code\_function}


% لیست موارد بدون شماره‌گذاری:
%\شروع{فقرات}
%    \فقره مورد اول 
%        \شروع{فقرات}
%            \فقره مورد اول داخلی 
%                \شروع{فقرات}
%                    \فقره اولین مورد خیلی داخلی 
%                    \فقره دومین مورد خیلی داخلی 
%                \پایان{فقرات}
%            \فقره مورد دوم داخلی 
%        \پایان{فقرات}
%    \فقره مورد دوم 
%\پایان{فقرات}

% لیست موارد شماره‌گذاری شده:
%\شروع{شمارش}
%    \فقره یک مورد 
%    \فقره یک مورد دیگر 
%    \فقره باز هم یک مورد دیگر
%\پایان{شمارش}

% برای نوشتن یک فرمول چند خطی که یک شماره بخوره. برای حالت‌های دیگر به سایت زیر مراجعه کنید:
% https://www.overleaf.com/learn/latex/Aligning_equations_with_amsmath
%\begin{align} 
%    \begin{split}
%        (x+y)^3 &= (x+y)^2(x+y) \\
%                &= (x^2+2xy+y^2)(x+y) \\
%                &= (x^3+2x^2y+xy^2) + (x^2y+2xy^2+y^3) \\
%                &= x^3+3x^2y+3xy^2+y^3
%    \end{split}                                 
%\end{align}

% نوشتن ماتریس و آرایه دوبعدی
%\begin{align}
%    A = 
%    \begin{bmatrix}
%        A_{11} & A_{21} \\
%        A_{21} & A_{22}
%    \end{bmatrix}
%\end{align}

% اضافه کردن تصویر با شرح. برای اضافه کردن تصویر بدون شرح، فقط دستور درج‌تصویر را استفاده کنید.
%\شروع{شکل}[h!]
%\تنظیم‌ازوسط
%  \درج‌تصویر[پهنا=0.5\پهنای‌متن]{img.jpg}
%\شرح{شرح تصویر}
%\برچسب{تصویر:۱}
%\پایان{شکل}
%\پایان{زیرمسائل}

%کد زیر را ببینید.
%%\lstinputlisting{source_filename.py}
%\begin{latin}
%\begin{lstlisting}
%# Import the modules
%import sys
%import random
%
%ans = True
%
%while ans:
%question = raw_input("Ask the magic 8 ball a question: (press enter to quit) ")
%
%answers = random.randint(1,8)
%
%if question == "":
%sys.exit()
%
%elif answers == 1:
%print "It is certain"
%
%elif answers == 2:
%print "Outlook good"
%
%elif answers == 3:
%print "You may rely on it"
%
%elif answers == 4:
%print "Ask again later"
%
%elif answers == 5:
%print "Concentrate and ask again"
%
%elif answers == 6:
%print "Reply hazy, try again"
%
%elif answers == 7:
%print "My reply is no"
%
%elif answers == 8:
%print "My sources say no"
%\end{lstlisting}
%\end{latin}
